% concatenated from main.tex
\documentclass[11pt]{article}
%\usepackage[margin=1in]{geometry}
\usepackage{amsfonts}
\usepackage{amsthm}
\usepackage{amsmath}
\usepackage{amssymb}
\usepackage{algorithm}
\usepackage{algorithmic}
\usepackage{mathtools}
\usepackage{hyperref}
% \usepackage{cleveref}
\usepackage{color}

\newcommand{\FF}{\mathbb{F}}
\newcommand{\ZZ}{\mathbb{Z}}
\newcommand{\0}{\mathbf{0}}
\newcommand{\1}{\mathbf{1}}
\DeclareMathOperator{\Hull}{Hull}
\DeclareMathOperator{\rank}{rank}
\DeclareMathOperator{\wt}{wt}
\DeclareMathOperator{\Aut}{Aut}
\DeclareMathOperator{\disc}{disc}

\theoremstyle{plain}
\newtheorem{theorem}{Theorem}[section]
\newtheorem{lemma}[theorem]{Lemma}
\newtheorem{corollary}[theorem]{Corollary}
\newtheorem{proposition}[theorem]{Proposition}
\newtheorem{conjecture}[theorem]{Conjecture}
\theoremstyle{definition}
\newtheorem{definition}[theorem]{Definition}
\newtheorem{example}[theorem]{Example}
\theoremstyle{remark}
\newtheorem{remark}[theorem]{Remark}

% \crefname{theorem}{Theorem}{Theorems}
% \crefname{lemma}{Lemma}{Lemmas}
% \crefname{corollary}{Corollary}{Corollaries}
% \crefname{proposition}{Proposition}{Propositions}

\title{The Closure of LCD-to-GI Reductions via Generalized Inner Products}

\author{Keita~Ishizuka
\thanks{K. Ishizuka is with the Information Technology R\&D Center, Mitsubishi Electric Corporation, Kanagawa, Japan (e-mail: keitaishizuka1994@gmail.com).}}
\date{}

\begin{document}

\maketitle

\begin{abstract}
The Permutation Equivalence Problem (PEP) for linear codes is a fundamental problem in coding theory and cryptography. A recent reduction shows that PEP for Linear Complementary Dual (LCD) codes reduces to Graph Isomorphism~(GI) via orthogonal projectors, but is restricted to codes with trivial hull. We prove that this approach extends to bilinear forms $M = aI + bJ$, and that no other nondegenerate symmetric form yields a valid reduction. A code is reducible if and only if its hull dimension is at most~$1$ with an explicit condition on the hull vector; in characteristic~$2$, only LCD codes are reducible. This establishes the closure of the orthogonal projector method. We derive exact enumeration formulas via character sums over quadratic forms and provide a polynomial-time reduction algorithm.
\end{abstract}

\medskip
\noindent\textbf{Keywords:} Permutation Equivalence Problem, Graph Isomorphism, LCD codes, Enumeration

\section{Introduction}

The Permutation Equivalence Problem (PEP) for linear codes asks whether two given codes are identical up to a coordinate permutation.
In coding theory, PEP underlies the classification of linear codes, a long-standing problem for which Sendrier's Support Splitting Algorithm~\cite{sendrier2000finding} remains the standard computational tool. In cryptography, the conjectured hardness of PEP and its close variants underlies a growing family of code-based signature schemes. The NIST PQC candidate LESS~\cite{LESS} originally relied on PEP and now uses the closely related Linear Equivalence Problem (LEP); the recent SPECK~\cite{SPECK} is based on the Permutation Equivalence of Codes and Kernels (PECK) problem, which combines features of PEP and the Permuted Kernel Problem (PKP) and reduces to PEP in the high-$q$ regime, where the scheme draws hard instances from codes with large hull. Understanding the structural and computational boundaries of PEP therefore has direct implications for both code classification and post-quantum cryptography.

The complexity of PEP is closely tied to that of Graph Isomorphism (GI). Petrank and Roth~\cite{petrank1997finding} showed that PEP is at least as hard as GI; conversely, Bardet et al.~\cite{bardet2019permutation} proved that PEP for Linear Complementary Dual (LCD) codes (those with trivial hull $\Hull(C) = C \cap C^\perp = \{0\}$) reduces to GI in polynomial time. Their construction maps each LCD code to a weighted graph via the orthogonal projector $\Pi_C = G^T(GG^T)^{-1}G$, with the property that two LCD codes are permutation-equivalent if and only if their projectors are isomorphic as weighted graphs. Combined with Babai's quasi-polynomial-time algorithm for GI~\cite{babai2016}, this yields a sub-exponential algorithm for PEP on LCD instances. However, the construction requires $GG^T$ to be nonsingular, ruling out codes with nontrivial hull. The natural question is: \emph{how far can the orthogonal projector method be pushed?}

\paragraph{Our contribution.}
We answer this question completely. First, we prove that the orthogonal projector $\Pi_{C,M} = MG^T(GMG^T)^{-1}G$ gives a valid reduction ($C_1 \cong C_2 \Leftrightarrow \Pi_{C_1,M} \cong \Pi_{C_2,M}$) if and only if $M = aI + bJ$ (Theorem~\ref{thm:main_characterization}). Second, we characterize exactly which codes are reducible: a code admits such a reduction if and only if its hull dimension is at most~$1$ with $(x,\1) \neq 0$ for the hull vector~$x$ (Theorem~\ref{thm:gi_reducible_count}); in characteristic~$2$, only LCD codes are reducible (Corollary~\ref{cor:char2}). Together, these results establish the closure of the orthogonal projector approach: codes with hull dimension $\geq 2$ are provably beyond its reach, regardless of the choice of nondegenerate symmetric bilinear form. We also derive exact enumeration formulas using Weil's theorem on quadratic forms (Theorem~\ref{thm:enumeration_main}) and provide a polynomial-time reduction algorithm (Section~\ref{sec:algorithm}).

\paragraph{Relation to prior work.}
A preliminary 3-page poster by the author~\cite{ishizuka2025ccsposter} introduced the family of permutation-invariant bilinear forms with structure matrix $aI + bJ$ and established the dimension bound $|\dim \Hull_M(C) - \dim \Hull(C)| \leq 1$. That work used these tools to accelerate the Support Splitting Algorithm~\cite{sendrier2000finding}. The present paper shares only this characterization and dimension bound with the poster (Lemma~\ref{lem:centralizer} and Theorem~\ref{thm:dimension_bound}), reproduced here for self-containment. The main results of this paper address the distinct question of GI-reducibility and are not present in~\cite{ishizuka2025ccsposter}.

\paragraph{Organization.}
Section~\ref{sec:preliminaries} reviews definitions and basic properties.
Section~\ref{sec:extension} introduces permutation-invariant bilinear forms and characterizes valid structure matrices.
Section~\ref{sec:characterization} characterizes GI-reducible codes and proves the closure result.
Section~\ref{sec:enumeration} derives enumeration formulas.
Section~\ref{sec:algorithm} presents the reduction algorithm.
Section~\ref{sec:conclusion} concludes with open questions.
Appendix~\ref{app:lcd_enumeration} collects standard LCD code enumeration formulas (Carlet et al.), and Appendix~\ref{app:type_epsilon} extends them to the type-$\epsilon$ setting needed in Section~\ref{sec:enumeration}.


\section{Preliminaries}
\label{sec:preliminaries}

Let $S_n$ denote the symmetric group on $n$ elements. For $\pi \in S_n$ and a vector $x = (x_1, \ldots, x_n) \in \FF_q^n$, we define $\pi(x) = (x_{\pi^{-1}(1)}, \ldots, x_{\pi^{-1}(n)})$. The corresponding permutation matrix $P_\pi \in \FF_q^{n \times n}$ is defined by $(P_\pi)_{ij} = 1$ if $j = \pi(i)$ and $0$ otherwise, so that $\pi(x) = xP_\pi$ for row vectors.

\subsection{Graphs}

In this paper, we consider directed graphs with self-loops and edge weights in $\FF_q$. A directed graph on $n$ vertices is represented by a matrix $A \in \FF_q^{n \times n}$ where $A_{ij}$ denotes the weight of the directed edge from vertex $i$ to vertex $j$. Self-loops are allowed ($A_{ii}$ may be nonzero). An \emph{undirected graph} is the special case where the adjacency matrix $A$ is symmetric.

Two graphs with adjacency matrices $A_1, A_2 \in \FF_q^{n \times n}$ are \emph{isomorphic} if there exists a permutation $\pi \in S_n$ such that $A_2 = P_\pi^T A_1 P_\pi$, where $P_\pi$ is the permutation matrix corresponding to $\pi$.

\begin{definition}[Graph Isomorphism Problem]
Given adjacency matrices $A_1$ and $A_2$ of two graphs, decide whether there exists a permutation $\pi \in S_n$ such that $A_2 = P_\pi^T A_1 P_\pi$.
\end{definition}

It is well known that directed graph isomorphism reduces to undirected graph isomorphism in polynomial time (by a standard gadget construction), and that edge weights can be eliminated in polynomial time by replacing each weighted edge with a gadget encoding the weight. Babai~\cite{babai2016} showed that GI can be solved in quasipolynomial time $2^{O(\log^c n)}$ for some constant $c$. Combined with the above reductions, directed graph isomorphism can also be solved in quasipolynomial time.

\subsection{Linear Codes}
Let $\FF_q$ denote the finite field of order $q$, where $q = p^m$ for some prime $p$ and positive integer $m$. The characteristic of $\FF_q$ is $p$. We denote by $\FF_q^n$ the $n$-dimensional vector space over $\FF_q$.
An $[n,k]_q$ code is a $k$-dimensional subspace of $\FF_q^n$.
The values $n$ and $k$ are called length and dimension, respectively.
A generator matrix of a code $C$ is a $k \times n$ matrix $G$ whose rows form a basis of $C$.
For vectors $v_1, \ldots, v_m \in \FF_q^n$, we denote by $\langle v_1, \ldots, v_m \rangle$ the subspace spanned by these vectors.
Two codes $C_1$ and $C_2$ of length $n$ over $\FF_q$ are \emph{permutation-equivalent} if there exists $\pi \in S_n$ such that $C_2 = C_1 P_\pi$, where $P_\pi$ is the permutation matrix corresponding to $\pi$.

\begin{definition}[Permutation Equivalence Problem]
Given two $[n,k]_q$ codes $C_1$ and $C_2$, decide whether there exists a permutation $\pi \in S_n$ such that $C_2 = C_1 P_\pi$.
\end{definition}

The standard inner product on $\FF_q^n$ is defined as $\langle x, y \rangle = \sum_{i=1}^n x_i y_i$ for $x = (x_1, \ldots, x_n)$ and $y = (y_1, \ldots, y_n)$ in $\FF_q^n$.
The \emph{dual code} of $C$ with respect to the standard inner product is $C^\perp = \{y \in \FF_q^n : \langle x, y \rangle = 0 \text{ for all } x \in C\}$.
If $C$ is an $[n,k]_q$ code, then $C^\perp$ is an $[n,n-k]_q$ code. Moreover, $(C^\perp)^\perp = C$.
The \emph{hull} of a code $C$ is defined as $\Hull(C) = C \cap C^\perp$.
A code $C$ is called linear complementary dual (LCD) code if $\Hull(C) = \{0\}$.
A matrix $\Pi \in \FF_q^{n \times n}$ is an \emph{orthogonal projector} onto a subspace $C \subseteq \FF_q^n$ (with respect to the standard inner product) if $\Pi^2 = \Pi$, $\text{Im}(\Pi) = C$, and for all $x \in \text{Im}(\Pi)$ and $y \in \ker(\Pi)$, we have $\langle x, y \rangle = 0$.
\begin{theorem}[Massey \cite{massey1992linear}]
\label{thm:massey}
Let $C$ be an $[n,k]$ linear code with generator matrix $G$. Then $C$ is LCD if and only if $GG^T$ is nonsingular. Moreover, if $C$ is LCD, then
$$\Pi_C = G^T(GG^T)^{-1}G$$
is the orthogonal projector from $\FF_q^n$ onto $C$.
\end{theorem}

\begin{theorem}[GI reduction, Bardet et al.\ \cite{bardet2019permutation}]
\label{thm:gi_reduction}
Let $C_1$ and $C_2$ be two $[n,k]$ linear codes with $\Hull(C_i) = \{0\}$ (i.e., LCD codes) for $i=1,2$. Then the permutation equivalence problem for $C_1$ and $C_2$ can be reduced in polynomial time to the graph isomorphism problem.
\end{theorem}

It is worthwhile to review the argument, since our main result extends it from the standard inner product to arbitrary symmetric bilinear forms. For an LCD code $C$ with generator matrix $G$, the projector $\Pi_C = G^T(GG^T)^{-1}G$ is symmetric, idempotent, and a complete invariant of $C$: one has $C = \mathrm{Im}(\Pi_C)$, and $\Pi_C$ depends only on $C$ (not on the choice of $G$). For any permutation matrix $P$, direct computation gives $\Pi_{CP} = P^T \Pi_C P$. Hence two LCD codes are permutation equivalent if and only if their projectors are conjugate by a permutation matrix, i.e., isomorphic as edge-weighted graphs on $n$ vertices whose adjacency matrices are $\Pi_{C_1}$ and $\Pi_{C_2}$; a standard gadget reduces weighted GI over $\FF_q$ to (unweighted) GI in polynomial time, completing the reduction.

The crux of this argument is the identity $\Pi_{CP} = P^T \Pi_C P$, which relies on $PIP^T = I$ (since $P$ is orthogonal). Section~\ref{sec:extension} replaces $I$ by a general symmetric structure matrix $M$ and asks: for which $M$ does the analogous identity $\Pi_{CP, M} = P^T \Pi_{C, M} P$ hold, so that the reduction goes through?

\subsection{Quadratic and Bilinear Forms}

We summarize the structural facts on quadratic and symmetric bilinear forms used throughout the paper; the standard references are~\cite{lidl1997finitefields} and~\cite{taylor1992geometry}. Throughout this subsection $q$ is odd unless stated otherwise.

A \emph{bilinear form} on $\FF_q^n$ is a map $B: \FF_q^n \times \FF_q^n \rightarrow \FF_q$ that is linear in each argument. Every bilinear form $B$ can be represented by a matrix $M \in \FF_q^{n \times n}$ such that $B(x, y) = xMy^T$ for all $x, y \in \FF_q^n$, where we view $x$ and $y$ as row vectors. We call $M$ the \emph{structure matrix} of $B$. A bilinear form is \emph{nondegenerate} if $\det(M) \neq 0$. Throughout this paper, we restrict to \emph{symmetric} bilinear forms (i.e., $M = M^T$): a non-symmetric structure matrix would necessitate distinguishing left and right dual codes, which falls outside our framework. We identify a bilinear form with its structure matrix with respect to the standard basis, and use the matrix notation $M$ directly instead of the function notation $B$.
The \emph{radical} of a symmetric bilinear form $B$ on a subspace $W \subseteq \FF_q^n$ is
\[
\mathrm{rad}(W) := \{v \in W : B(v, w) = 0 \text{ for all } w \in W\}.
\]
$B|_W$ is \emph{nondegenerate} if $\mathrm{rad}(W) = \{0\}$; equivalently, the Gram matrix of any basis of~$W$ is nonsingular. A vector $x$ is \emph{isotropic} if  $B(x,x) = 0$.
The \emph{discriminant} of a nondegenerate $B$ on $\FF_q^n$ is the class $\disc(B) := \det(M) \cdot (\FF_q^*)^2 \in \FF_q^*/(\FF_q^*)^2$. This is invariant under change of basis $M \mapsto QMQ^T$ ($Q \in \mathrm{GL}_n$, $\det(QMQ^T) = \det(Q)^2 \det(M)$). Two nondegenerate symmetric bilinear forms on $\FF_q^n$ ($q$ odd) are \emph{isometric} (related by such a change of basis) if and only if they have the same discriminant.
A two-dimensional nondegenerate $B$ with isotropic basis $\{u, v\}$ satisfying $B(u,v) = 1$ is called a \emph{hyperbolic plane}; its Gram matrix is $\begin{psmallmatrix} 0 & 1 \\ 1 & 0 \end{psmallmatrix}$, so $\disc = -1$. A subspace on which $B$ vanishes only at the origin is called \emph{anisotropic}.

The following three structural results from the theory of bilinear forms over $\FF_q$ are used repeatedly; all are standard~\cite{taylor1992geometry}.

\begin{lemma}[Hyperbolic completion]\label{lem:hyperbolic_completion}
Let $q$ be odd and let $B$ be a nondegenerate symmetric bilinear form on $V = \FF_q^n$. For every nonzero isotropic vector $u \in V$, there exists $v \in V$ such that $\{u, v\}$ spans a hyperbolic plane.
\end{lemma}

\begin{lemma}[Existence of an orthogonal basis]\label{lem:orth_basis}
Let $q$ be odd. Every nondegenerate symmetric bilinear form $B$ on $\FF_q^n$ admits an orthogonal basis: the structure matrix $M$ is congruent to $\mathrm{diag}(d_1, \ldots, d_n)$ for some $d_1, \ldots, d_n \in \FF_q^*$.
\end{lemma}

\begin{theorem}[Witt decomposition]\label{thm:witt_decomp}
Let $q$ be odd and let $B$ be a nondegenerate symmetric bilinear form on $V = \FF_q^n$. Then $V$ decomposes as
\[
V = H \perp W,
\]
where $H$ is an orthogonal sum of hyperbolic planes (the \emph{hyperbolic part}) and $W$ is anisotropic with $\dim W \leq 2$. The decomposition is unique up to isometry.
\end{theorem}

\begin{theorem}[Witt's extension theorem]\label{thm:witt_extension}
Let $q$ be odd and let $B$ be a nondegenerate symmetric bilinear form on $V = \FF_q^n$. Any isometry $\phi: W_1 \to W_2$ between subspaces $W_1, W_2 \subseteq V$ extends to an isometry $\Phi: V \to V$ lying in the orthogonal group $\mathbb{O}(V, B)$.
\end{theorem}

A \emph{quadratic form} $Q(x) = \sum_{i,j} Q_{ij} x_i x_j$ over $\FF_q^n$ is a homogeneous polynomial of degree~$2$. Its \emph{associated (symmetric) bilinear form} is $B_Q(x,y) = Q(x+y) - Q(x) - Q(y)$, satisfying $B_Q(x,x) = 2Q(x)$. For $q$ odd, polarization gives a bijection between quadratic forms and symmetric bilinear forms via $Q \mapsto B_Q$ and $B \mapsto Q(x) := \tfrac{1}{2}B(x,x)$. We henceforth identify a symmetric bilinear form with its symmetric matrix~$M$ (so $B(x,y) = xMy^T$ for row vectors) and a quadratic form with its symmetric Gram matrix.
When $q$ is odd, the \emph{quadratic character} $\chi: \FF_q^* \to \{-1, +1\}$ is defined by $\chi(a) = 1$ if $a \in (\FF_q^*)^2$ and $\chi(a) = -1$ otherwise. We extend $\chi$ to $\FF_q$ by setting $\chi(0) = 0$. The discriminant class of $B$ is encoded by $\chi(\det M) \in \{+1, -1\}$. The following classical result counts solutions to quadratic forms.

\begin{theorem}[Weil; Lidl and Niederreiter {\cite{lidl1997finitefields}}]\label{thm:weil_lidl_niederreiter}
Let $Q(x) = \sum_{i,j} Q_{ij} x_i x_j$ be a nondegenerate quadratic form over $\mathbb{F}_q^n$ where $q$ is odd, and let $N(Q) = |\{x \in \mathbb{F}_q^n : Q(x) = 0\}|$ be the number of its solutions. Then
$$
N(Q) = \begin{cases}
    q^{n-1} & \textrm{ if $n$ is odd,} \\
    q^{n-1} + \chi((-1)^{n/2} \det(Q))(q-1)q^{(n-2)/2} & \textrm{ if $n$ is even,}
\end{cases}
$$
where $\chi$ is the quadratic character of $\mathbb{F}_q$ and $\det(Q)$ denotes the determinant of the matrix of $Q$.
\end{theorem}

\section{Permutation-Invariant Bilinear Forms}
\label{sec:extension}

The reduction of Bardet et al.\ (Theorem~\ref{thm:gi_reduction}) applies only to LCD codes, those with trivial hull $\Hull(C) = \{0\}$. For non-LCD codes, the Gram matrix $GG^T$ is singular, making the orthogonal projector undefined. To extend the reduction to codes with nonzero hull dimension, we replace the standard inner product with more general bilinear forms parametrized by matrices $M$. This yields a generalized notion of $M$-LCD codes and corresponding orthogonal projectors. The key question: which $M$ preserve the reduction structure?

\subsection{Generalized Massey Theorem}

Fix a symmetric structure matrix $M \in \FF_q^{n \times n}$. For a code $C \subseteq \FF_q^n$, we define its \emph{$M$-dual code} by
\[
C^{\perp_M} := \{y \in \FF_q^n : xMy^T = 0 \text{ for all } x \in C\},
\]
its \emph{$M$-hull} by $\Hull_M(C) := C \cap C^{\perp_M}$, and call $C$ an \emph{$M$-LCD code} if $\Hull_M(C) = \{0\}$. Setting $M = I$ recovers the standard inner-product notions of dual code, hull, and LCD code. With this generalized duality in place, Massey's theorem extends to arbitrary nondegenerate symmetric bilinear forms as follows.

\begin{lemma}
\label{lem:hull_dimension}
Let $C$ be an $[n,k]$ code with generator matrix $G$, and let $M$ be a structure matrix. Then $\dim \Hull_M(C) = k - \rank(GMG^T)$.
\end{lemma}
\begin{proof}
$\Hull_M(C) = C \cap C^{\perp_M} = \{cG : c \in \FF_q^k,\ cGMG^T = 0\}$. The map $c \mapsto cG$ is injective on $\FF_q^k$ (since $G$ has full row rank), so $\dim \Hull_M(C) = \dim\{c \in \FF_q^k : cGMG^T = 0\} = k - \rank(GMG^T)$.
\end{proof}

\begin{theorem}[Generalized Massey Theorem]
\label{thm:generalized_massey}
Let $C$ be an $[n,k]$ linear code with generator matrix $G$, and let $M$ be a nondegenerate symmetric structure matrix. Then $\Hull_M(C) = \{0\}$ if and only if $GMG^T$ is nonsingular. Moreover, if $\Hull_M(C) = \{0\}$, then
$$\Pi_{C,M} = MG^T(GMG^T)^{-1}G$$
is the orthogonal projector from $\FF_q^n$ onto $C$ with respect to $M$, satisfying $\Pi_{C,M}^2 = \Pi_{C,M}$, $\text{Im}(\Pi_{C,M}) = C$, and for all $v, w \in \FF_q^n$ with $v\Pi_{C,M} = v$ and $w\Pi_{C,M} = 0$, we have $vMw^T = 0$.
\end{theorem}

\begin{proof}
By Lemma~\ref{lem:hull_dimension}, $\dim \Hull_M(C) = k - \rank(GMG^T)$. Thus $\Hull_M(C) = \{0\}$ if and only if $\rank(GMG^T) = k$, i.e., $GMG^T$ is nonsingular.
For the projector formula, we verify that $\Pi_{C,M} = MG^T(GMG^T)^{-1}G$ satisfies the three defining properties. Idempotence $\Pi_{C,M}^2 = \Pi_{C,M}$ follows by direct computation.

For the image, let $v = cG \in C$ where $c \in \mathbb{F}_q^k$.
Then
\begin{align*}
v\Pi_{C,M}
&= cG \cdot MG^T(GMG^T)^{-1}G\\
&= c(GMG^T)(GMG^T)^{-1}G\\
&= cG = v,
\end{align*}
so $v$ lies in the image of $\Pi_{C,M}$. Conversely, any vector in the image has the form $w\Pi_{C,M} = wMG^T(GMG^T)^{-1}G$ for some row vector $w \in \mathbb{F}_q^n$; setting $c = wMG^T(GMG^T)^{-1} \in \mathbb{F}_q^k$ gives $w\Pi_{C,M} = cG \in C$. Hence the image equals $C$.

For orthogonality, let $v, w \in \mathbb{F}_q^n$ satisfy $v\Pi_{C,M} = v$ and $w\Pi_{C,M} = 0$. Then $v = cG$ for some $c \in \mathbb{F}_q^k$, and from $wMG^T(GMG^T)^{-1}G = 0$ we obtain $wMG^T = 0$ since $(GMG^T)^{-1}$ is invertible and $G$ has full row rank. Therefore $vMw^T = cGMw^T = c(wMG^T)^T = 0$, as required.
\end{proof}

\subsection{Characterization and Construction}

With Generalized Massey in hand, we can extend Bardet et al.'s reduction to $M$-LCD codes for arbitrary structure matrices $M$. However, a crucial question arises: Which choices of $M$ yield a valid reduction to graph isomorphism? That is, for which $M$ does the orthogonal projector map $C \mapsto \Pi_{C,M}$ satisfy
\[
C_1 \cong C_2 \iff \Pi_{C_1,M} \cong \Pi_{C_2,M}
\]
for all $M$-LCD codes $C_1, C_2$?
To answer this, we first collect some auxiliary lemmas.

\begin{lemma}[Centralizer of the permutation action, \cite{ishizuka2025ccsposter}]
\label{lem:centralizer}
Let $\FF$ be any field and let $S_n$ act on $\FF^n$ by permuting coordinates. Then the centralizer of this action is
\[
\{X \in \FF^{n \times n} : P^T X P = X \text{ for all } P \in S_n\} = \operatorname{span}\{I, J\},
\]
where $J$ is the all-ones matrix.
\end{lemma}

\begin{proof}
The symmetric group $S_n$ acts on the index set $\{1, \ldots, n\} \times \{1, \ldots, n\}$ by $\pi \cdot (i,j) = (\pi(i), \pi(j))$. This action has exactly two orbits: the diagonal $\Delta = \{(i,i) : i = 1, \ldots, n\}$ and the off-diagonal $\Omega = \{(i,j) : i \neq j\}$. Since $X$ is invariant under all permutations, its entries must be constant on each orbit. Thus $X_{ii} = a$ for all $i$ (diagonal entries) and $X_{ij} = b$ for all $i \neq j$ (off-diagonal entries), giving $X = aI + bJ$ for some $a, b \in \FF$.
\end{proof}

\begin{lemma}[Proportionality of quadratic forms]
\label{lem:quadratic_proportionality}
Let $q$ be odd and $n \geq 3$. If $Q_1, Q_2$ are nondegenerate quadratic forms on $\FF_q^n$ with the same zero set, then $Q_2 = \lambda Q_1$ for some $\lambda \in \FF_q^*$.
\end{lemma}
\begin{proof}
Pick $u_0 \in \FF_q^n$ with $Q_1(u_0) \neq 0$ (and hence $Q_2(u_0) \neq 0$, by the equal-zero-set hypothesis); set $\lambda := Q_2(u_0)/Q_1(u_0) \in \FF_q^*$. For any nonzero isotropic vector~$w$, the line $\{u_0 + t w : t \in \FF_q\}$ yields polynomials
\[
f_i(t) := Q_i(u_0 + t w) = Q_i(u_0) + 2 t B_i(u_0, w),
\]
where $B_i$ denotes the symmetric bilinear form polarized from $Q_i$. Both $f_1, f_2$ are linear in~$t$, with the same zero set in~$\FF_q$. Equating root structures forces $B_2(u_0, w) = \lambda B_1(u_0, w)$: if $B_1(u_0, w) = 0$ then $f_1 \equiv Q_1(u_0)$ has empty zero set, so $B_2(u_0, w) = 0$ as well; otherwise $f_1$ has the unique root $-Q_1(u_0)/(2 B_1(u_0, w))$, and equating with the unique root of $f_2$ gives the desired identity.

For $n \geq 3$ and $q$ odd, the isotropic vectors of any nondegenerate quadratic form $Q$ span $\FF_q^n$. By the Witt decomposition (Theorem~\ref{thm:witt_decomp}), $\FF_q^n = H \perp W$ with $H$ an orthogonal sum of hyperbolic planes and $W$ anisotropic of dimension $\leq 2$: the hyperbolic bases of $H$ are isotropic and span~$H$; for any $w \in W \setminus \{0\}$, surjectivity of the form on a hyperbolic plane provides $h \in H$ with $Q(h) = -Q(w)$, so $h + w$ is isotropic with nonzero $W$-component. Combining these vectors yields a spanning set of $\FF_q^n$ consisting of isotropic vectors. Hence the linear functional $v \mapsto B_2(u_0, v) - \lambda B_1(u_0, v)$ vanishes on a spanning set, so $B_2(u_0, \cdot) = \lambda B_1(u_0, \cdot)$ identically.

Repeating the argument with any $u_0' \in \FF_q^n$ satisfying $Q_1(u_0') \neq 0$ gives $B_2(u_0', \cdot) = \lambda_{u_0'} B_1(u_0', \cdot)$ for some $\lambda_{u_0'} \in \FF_q^*$. Choosing $u_0'$ so that $B_1(u_0, u_0') \neq 0$ (possible by non-degeneracy of $B_1$ together with $\dim\{v : B_1(u_0, v) \neq 0\} = n - 1$), the symmetry of the bilinear forms yields
\[
\lambda B_1(u_0, u_0') = B_2(u_0, u_0') = B_2(u_0', u_0) = \lambda_{u_0'} B_1(u_0', u_0) = \lambda_{u_0'} B_1(u_0, u_0'),
\]
so $\lambda_{u_0'} = \lambda$. Since $\{u : Q_1(u) \neq 0\}$ spans $\FF_q^n$, $B_2 = \lambda B_1$ identically, hence $Q_2 = \lambda Q_1$.
\end{proof}

\begin{theorem}[Characterization of valid reductions]
\label{thm:main_characterization}
Let $M$ be a nondegenerate symmetric structure matrix and $n \geq 3$. Then
\[
C_1 \cong C_2 \iff \Pi_{C_1,M} \cong \Pi_{C_2,M}
\]
holds for all $M$-LCD codes $C_1, C_2$ if and only if $M = aI + bJ$ for some $a, b \in \FF_q$ with $a \neq 0$ and $a + nb \neq 0$.
\end{theorem}

\begin{proof}
$(\Leftarrow)$
Suppose $M = aI + bJ$. Then $PMP^T = M$ for all permutation matrices~$P$ (since $PIP^T = I$ and $PJP^T = J$). For the forward implication of the biconditional, if $\Pi_{C_2,M} = P^T\Pi_{C_1,M}P$ for some permutation~$P$, then $C_2 = \mathrm{Im}(\Pi_{C_2,M}) = C_1P$ by the image property $\mathrm{Im}(\Pi_{C,M}) = C$ (Theorem~\ref{thm:generalized_massey}). This holds for any nondegenerate~$M$. For the reverse implication, suppose $C_2 = C_1 P$, so that $G_2 := G_1 P$ is a generator matrix of $C_2$. Since $M = aI + bJ$ and every permutation matrix $P$ satisfies $P\1 = \1$, we have $PI = IP = P$ and $PJ = JP = J$, hence $PM = MP$, and combined with $PP^T = I$, also $PMP^T = M$. Taking the transpose of $PM = MP$ together with $M = M^T$ yields $MP^T = P^T M$. Therefore,
\begin{align*}
\Pi_{C_2,M}
&= M G_2^T \bigl(G_2 M G_2^T\bigr)^{-1} G_2 \\
&= M (G_1 P)^T \bigl(G_1 P M (G_1 P)^T\bigr)^{-1} G_1 P \\
&= M P^T G_1^T \bigl(G_1 (PMP^T) G_1^T\bigr)^{-1} G_1 P \\
&= M P^T G_1^T (G_1 M G_1^T)^{-1} G_1 P
   && \text{(using $PMP^T = M$)} \\
&= P^T M G_1^T (G_1 M G_1^T)^{-1} G_1 P
   && \text{(using $MP^T = P^T M$)} \\
&= P^T \Pi_{C_1,M} P.
\end{align*}

$(\Rightarrow)$
Suppose the biconditional holds. We first observe that for any $M$-LCD code~$C$ and any permutation~$P$, the code $CP$ must also be $M$-LCD: indeed, $C \cong CP$ trivially, so the forward direction of the biconditional applied with $C_1 = C$, $C_2 = CP$ requires $\Pi_{CP,M}$ to be defined, which by Theorem~\ref{thm:generalized_massey} forces $CP$ to be $M$-LCD. In particular, applying this to the one-dimensional $M$-LCD codes $\langle u \rangle$ (those with $uMu^T \neq 0$) yields $uMu^T = 0 \Leftrightarrow u(PMP^T)u^T = 0$ for all $u \in \FF_q^n$ and all~$P$, so the quadratic forms $u \mapsto uMu^T$ and $u \mapsto u(PMP^T)u^T$ share the same zero set. Both are nondegenerate (since $M$ is nondegenerate symmetric), so by Lemma~\ref{lem:quadratic_proportionality}, $PMP^T = \lambda_P M$ for some $\lambda_P \in \FF_q^*$.

The map $P \mapsto \lambda_P$ is a group homomorphism $S_n \to \FF_q^*$. For a transposition $\tau$, $\tau^2 = e$ gives $\lambda_\tau^2 = 1$, so $\lambda_\tau \in \{\pm 1\}$. If $\lambda_\tau = -1$ for some $\tau = (i\;j)$, then for $k, l \notin \{i, j\}$ (which $\tau$ fixes), $M_{kl} = (\tau M \tau^T)_{kl} = -M_{kl}$, forcing $M_{kl} = 0$ in odd characteristic; ranging over all transpositions yields $M = 0$, contradicting non-degeneracy. Hence $\lambda_\tau = 1$ for every transposition. Since transpositions generate $S_n$, $\lambda_P = 1$ for all $P \in S_n$, i.e., $PMP^T = M$, and Lemma~\ref{lem:centralizer} gives $M = aI + bJ$.
\end{proof}

Based on the above theorem, we say a code $C$ is \emph{GI-reducible} (via the orthogonal projector method) if there exists a nondegenerate symmetric structure matrix $M=aI+bJ$ such that every code in the permutation class of $C$ is $M$-LCD and the biconditional
\[
C_1 \cong C_2 \iff \Pi_{C_1,M} \cong \Pi_{C_2,M}
\]
holds for all $C_1, C_2$ in that class. The theorem of Bardet et al.\ shows that all LCD codes are GI-reducible (with $M = I$).

\begin{remark}[Non-degeneracy of $M = aI + bJ$]
\label{rem:nondegeneracy}
A direct computation shows
\[
\det(aI + bJ) = a^{n-1}(a + nb),
\]
where $n$ is treated as an element of $\mathbb{F}_q$. Thus $M = aI + bJ$ is nondegenerate if and only if $a \neq 0$ and $a + nb \neq 0$. (When $p \mid n$, the second condition reduces to $a \neq 0$.) Throughout this paper, when we write $M = aI + bJ$, we implicitly assume these conditions hold.
\end{remark}

\begin{example}
\label{ex:gi_reduction}
Consider the $[4, 2]_3$ code $C$ over $\FF_3$ with generator matrix $G = \begin{psmallmatrix} 1 & 1 & 0 & 0 \\ 0 & 1 & 1 & 0 \end{psmallmatrix}$, and let $P \in S_4$ be the permutation matrix that swaps coordinates~$1$ and~$2$. Define $C' := CP$, with generator matrix $G' = GP = \begin{psmallmatrix} 1 & 1 & 0 & 0 \\ 1 & 0 & 1 & 0 \end{psmallmatrix}$.
Under the standard inner product, $GG^T = G'(G')^T = \begin{psmallmatrix} 2 & 1 \\ 1 & 2 \end{psmallmatrix}$ has $\det = 0$ in $\FF_3$, so neither $C$ nor $C'$ is LCD: the projectors $\Pi_C, \Pi_{C'}$ are undefined and the reduction of Bardet et al.\ does not apply.

With $M = I + J$, however, both codes are $M$-LCD: a direct computation gives $GMG^T = G'M(G')^T = \begin{psmallmatrix} 0 & 2 \\ 2 & 0 \end{psmallmatrix}$, with $\det = 2 \neq 0$ in $\FF_3$. The generalized orthogonal projectors $\Pi_{C,M} = MG^T(GMG^T)^{-1}G$ and $\Pi_{C',M} = M(G')^T(G'M(G')^T)^{-1}G'$ evaluate to
\[
\Pi_{C,M} = \begin{pmatrix}
1 & 1 & 0 & 0 \\
0 & 0 & 0 & 0 \\
0 & 1 & 1 & 0 \\
1 & 2 & 1 & 0
\end{pmatrix},
\qquad
\Pi_{C',M} = \begin{pmatrix}
0 & 0 & 0 & 0 \\
1 & 1 & 0 & 0 \\
1 & 0 & 1 & 0 \\
2 & 1 & 1 & 0
\end{pmatrix},
\]
and one verifies $\Pi_{C',M} = P^T \Pi_{C,M} P$ directly (swapping rows and columns~$1$ and~$2$). Reading $\Pi_{C,M}$ and $\Pi_{C',M}$ as weighted-adjacency matrices on $4$ vertices, the codes $C$ and $C'$ are recognised as permutation equivalent through the graph isomorphism witnessed by $P$, even though Bardet's standard-inner-product reduction fails.
\end{example}

\section{Complete Characterization}
\label{sec:characterization}

Having established the construction method, we now characterize which codes are GI-reducible. We will show in Theorem~\ref{thm:dimension_bound} that any $M$-LCD code (for $M = aI + bJ$) is either LCD with respect to the standard inner product, or has $\dim(\Hull(C)) \leq 1$.

\begin{lemma}[Rank Perturbation]
\label{lemma:rank_perturbation}
For any matrices $A, B \in \FF_q^{n \times n}$, we have
$|\rank(A + B) - \rank(A)| \leq \rank(B)$.
\end{lemma}
\begin{proof}
Since $\text{Im}(A+B) \subseteq \text{Im}(A) + \text{Im}(B)$, we have $\rank(A+B) \leq \rank(A) + \rank(B)$. Applying this to $A = (A+B) + (-B)$ gives $\rank(A) \leq \rank(A+B) + \rank(B)$. Combining yields the result.
\end{proof}

\begin{theorem}[Dimension Bound, \cite{ishizuka2025ccsposter}]
\label{thm:dimension_bound}
Let $C$ be an $[n, k]$ code over $\FF_q$, and let $B$ be a bilinear form with nondegenerate structure matrix $M = aI + bJ$. Then
\[
|\dim \Hull_M(C) - \dim \Hull(C)| \leq 1.
\]
\end{theorem}

\begin{proof}
Let $G$ be a generator matrix of $C$. By Lemma~\ref{lem:hull_dimension}, $\dim \Hull_M(C) = k - \rank(GMG^T)$ and $\dim \Hull(C) = k - \rank(GG^T)$. Since $GMG^T = G(aI + bJ)G^T = aGG^T + bvv^T$ where $v = G\1^T$, we have $\rank(GMG^T) = \rank(GG^T + \frac{b}{a}vv^T)$. Thus $|\dim \Hull_M(C) - \dim \Hull(C)| = |\rank(GMG^T) - \rank(GG^T)| \leq \rank(vv^T) \leq 1$ by Lemma~\ref{lemma:rank_perturbation}.
\end{proof}

Now, we establish a complete characterization of $M$-LCD codes for $M=aI+bJ$. The following lemmas are standard results, whose proofs we include for completeness.

\begin{lemma}[Orthogonal decomposition]
\label{lem:orthogonal_decomposition}
Any $[n,k]_q$ code $C$ admits a decomposition
$$
C = \Hull(C) \perp U,
$$
where $U$ is any subspace of $C$ complementary to $\Hull(C)$; every such $U$ is an LCD code.
\end{lemma}
\begin{proof}
Let $H = \Hull(C) = C \cap C^\perp$ and let $U$ be any complement of $H$ in $C$. By the definition of hull, $C = H \perp U$.
Suppose there exists nonzero $x \in \Hull(U)$.
Then $x \in \Hull(C)$, contradicting $\Hull(C) \cap U = \{0\}$.
Therefore $U$ is LCD.
\end{proof}

\begin{lemma}[Schur complement]
\label{lem:schur_complement}
Let $A \in \FF_q^{k \times k}$ be invertible, $B \in \FF_q^{k \times m}$, $C \in \FF_q^{m \times k}$, and $D \in \FF_q^{m \times m}$. Then
\[
\det \begin{psmallmatrix} A & B \\ C & D \end{psmallmatrix} = \det(A) \cdot \det(D - CA^{-1}B).
\]
\end{lemma}
\begin{proof}
$\begin{psmallmatrix} I & 0 \\ -CA^{-1} & I \end{psmallmatrix} \begin{psmallmatrix} A & B \\ C & D \end{psmallmatrix} = \begin{psmallmatrix} A & B \\ 0 & D - CA^{-1}B \end{psmallmatrix}$. Taking determinants gives the result.
\end{proof}

\begin{theorem}[Complete Characterization]
\label{thm:gi_reducible_count}
An $[n,k]_q$ code $C$ is GI-reducible if and only if $C$ is LCD, or $\dim(\Hull(C)) = 1$ and $x \in \Hull(C)$ satisfies $(x,\1) \neq 0$.
\end{theorem}

\begin{proof}
LCD codes are always GI-reducible by the reduction of Bardet et al.\ (Theorem~\ref{thm:gi_reduction}). For codes with $\dim(\Hull(C)) = 1$, by Lemma~\ref{lem:orthogonal_decomposition}, we have $C = \langle x \rangle \perp U$ where $\Hull(C) = \langle x \rangle$ and $U$ is an LCD code of dimension $k-1$. We show that $C$ is GI-reducible if and only if $(x,\1) \neq 0$.

Define $G = \begin{psmallmatrix} x \\ G' \end{psmallmatrix}$ where $G'$ generates $U$, and let $w = G'\1^T$. Since $x \in \Hull(C)$, we have $(x,x) = 0$ and $xG'^T = 0$, so
\[
GMG^T = \begin{pmatrix} b(x,\1)^2 & b(x,\1) w^T \\ b(x,\1) w & aG'G'^T + bww^T \end{pmatrix}.
\]
If $(x,\1) = 0$, then the first row and column vanish, so $GMG^T$ is singular for all $(a,b)$: $C$ is not $M$-LCD for any $M = aI + bJ$.

If $(x,\1) \neq 0$ and $b \neq 0$, Lemma~\ref{lem:schur_complement} with $A = b(x,\1)^2$ gives
\begin{align*}
    \det(GMG^T)
    &= b(x,\1)^2 \cdot \det(aG'G'^T + bww^T - bww^T)\\
    &= b(x,\1)^2 \cdot a^{k-1} \cdot \det(G'G'^T).
\end{align*}
Since $U$ is LCD, $\det(G'G'^T) \neq 0$, so $\det(GMG^T) \neq 0$ for any $a \neq 0$, $b \neq 0$ with $a + nb \neq 0$.

Finally, by Theorem~\ref{thm:dimension_bound}, codes with $\dim(\Hull(C)) \geq 2$ cannot be $M$-LCD for any structure matrix $M = aI + bJ$, since $\dim(\Hull_M(C)) \geq \dim(\Hull(C)) - 1 \geq 1$ implies $C$ is not $M$-LCD. Therefore, such codes are not GI-reducible.
\end{proof}

\begin{corollary}\label{cor:char2}
If $q = 2^m$, then $C$ is GI-reducible if and only if $C$ is LCD.
\end{corollary}
\begin{proof}
In characteristic $2$, $(x,x) = \sum x_i^2 = (\sum x_i)^2 = (x,\1)^2$, so $x \in \Hull(C)$ implies $(x,x) = 0$ and hence $(x,\1) = 0$. By Theorem~\ref{thm:gi_reducible_count}, no code with $\dim(\Hull(C)) = 1$ is GI-reducible.
\end{proof}

By Corollary~\ref{cor:char2}, the GI-reducible codes over $\FF_{2^m}$ are precisely the LCD codes, whose enumeration is given in~\cite{carlet2019new,lishiling2025} (see Appendix~\ref{app:lcd_enumeration}). The enumeration in Section~\ref{sec:enumeration} therefore assumes $q$ is odd; the algorithm of Section~\ref{sec:algorithm} is stated for general $q$.

\section{Counting GI-Reducible Codes}
\label{sec:enumeration}

Having characterized which codes are GI-reducible (Theorem~\ref{thm:gi_reducible_count}), we now enumerate them. By the orthogonal decomposition (Lemma~\ref{lem:orthogonal_decomposition}), the count splits into LCD codes and hull dimension 1 codes with $(x,\1) \neq 0$. Since formulas for $L(n,k,q)$ counting LCD codes are known~\cite{carlet2019new,lishiling2025}, the key is computing $K(n,q) = |\{x \in \FF_q^n : (x,x) = 0, (x,\1) \neq 0\}|$. We compute this using Weil's theorem on quadratic forms (Theorem~\ref{thm:weil_lidl_niederreiter}).

\begin{lemma}[Isotropic Vectors with Nonzero Coordinate Sum]
\label{lem:quadratic_count}
Let $q$ be odd with characteristic $p$, and let $n \geq 2$ be a positive integer.
The number $K(n,q)$ of vectors $x \in \mathbb{F}_q^n$ with $(x,x) = 0$ and $(x, \1) \neq 0$ is given by
$$
K(n,q) = \begin{cases}
q^{n-2}(q-1), & p \mid n, \\
q^{n-2}(q-1) - \chi((-1)^{(n-1)/2} n)(q-1)q^{(n-3)/2}, & p \nmid n,\ n \text{ odd,} \\
q^{n-2}(q-1) + \chi((-1)^{n/2})(q-1)q^{(n-2)/2}, & p \nmid n,\ n \text{ even,}
\end{cases}
$$
where $\chi$ is the quadratic character of $\mathbb{F}_q$.
\end{lemma}

\begin{proof}
We compute $K(n,q) = N(n,q) - M(n,q)$, where $N(n,q) = |\{x \in \FF_q^n : (x,x) = 0\}|$ and $M(n,q) = |\{x \in \FF_q^n : (x,x) = 0, (x,\1) = 0\}|$.

Applying Theorem~\ref{thm:weil_lidl_niederreiter} to $Q(x) = \sum_{i=1}^n x_i^2$ (which is nondegenerate with $\det(Q) = 1$) gives
$$
N(n,q) = \begin{cases}
q^{n-1} & \text{if } n \text{ is odd,} \\
q^{n-1} + \chi((-1)^{n/2})(q-1)q^{(n-2)/2} & \text{if } n \text{ is even.}
\end{cases}
$$

For $M(n,q)$, we count solutions on the hyperplane $H = \{x : (x,\1) = 0\}$ of dimension $n-1$. Since $H^\perp = \langle \1 \rangle$ and $(\1,\1) = n$, the radical of $Q|_H$ is $H \cap H^\perp = H \cap \langle \1 \rangle$, which is $\langle \1 \rangle$ if $p \mid n$ and $\{0\}$ if $p \nmid n$.

\emph{Case $p \nmid n$}: $Q|_H$ is nondegenerate of dimension $n-1$. To apply Weil's theorem, we need $\det(Q|_H)$. Since $\FF_q^n = H \perp \langle \1 \rangle$ and $Q(\1) = n$, we have $\det(Q) = \det(Q|_H) \cdot n \cdot d^2$ for some nonzero $d$, so $\chi(\det(Q|_H)) = \chi(\det(Q)) \cdot \chi(n)^{-1} = \chi(n)$. Thus
$$
M(n,q) = \begin{cases}
q^{n-2} & \text{if } n \text{ is even,} \\
q^{n-2} + \chi((-1)^{(n-1)/2} n)(q-1)q^{(n-3)/2} & \text{if } n \text{ is odd.}
\end{cases}
$$

\emph{Case $p \mid n$}: $\1 \in H$ and $Q|_H$ has radical $\langle \1 \rangle$. The quotient form $\bar{Q}$ on $\bar{H} = H/\langle \1 \rangle$ is nondegenerate of dimension $n-2$. Since $Q(x + t\1) = Q(x)$ for all $x \in H$ and $t \in \FF_q$, each zero of $\bar{Q}$ lifts to $q$ zeros in $H$. In the basis $\{\bar{g}_i = \overline{e_i - e_n} : i = 1, \ldots, n-2\}$ of $\bar{H}$, the matrix of $\bar{Q}$ is $I_{n-2} + J_{n-2}$ (diagonal entries~$2$, off-diagonal entries~$1$). By the determinant formula in Remark~\ref{rem:nondegeneracy} applied with $a = b = 1$ and dimension $n-2$, $\det(I_{n-2} + J_{n-2}) = 1^{n-3}(1 + (n-2)) = n - 1 = -1$ in $\FF_q$ (using $p \mid n$). Applying Weil's theorem and multiplying by $q$:
$$
M(n,q) = \begin{cases}
q^{n-2} & \text{if } n \text{ is odd,} \\
q^{n-2} + \chi((-1)^{n/2})(q-1)q^{(n-2)/2} & \text{if } n \text{ is even,}
\end{cases}
$$
where we used $\chi((-1)^{(n-2)/2} \cdot (-1)) = \chi((-1)^{n/2})$.

Subtracting $M(n,q)$ from $N(n,q)$: when $p \mid n$, $K = N - M = q^{n-2}(q-1)$ regardless of parity (in the even case, the $\chi$ terms cancel). The remaining cases yield the stated formula.
\end{proof}

To count hull-dimension-1 codes, we fix a hull line $\langle x \rangle$ and enumerate codes with that hull. The structure of the residual bilinear form on $V = x^\perp / \langle x \rangle$ governs this count.

\begin{lemma}[Hull-Line Counting]\label{lem:hull_line_count}
Let $x \in \FF_q^n$ satisfy $(x,x) = 0$ and $(x,\1) \neq 0$. The number of $[n,k]_q$ codes $C$ with $\Hull(C) = \langle x \rangle$ equals $L^{\chi(-1)}(n-2,\, k-1,\, q)$, where $L^\epsilon(n,k,q)$ for $\epsilon \in \{+1, -1\}$ denotes the number of $k$-dimensional subspaces of $\FF_q^n$ that are LCD with respect to a nondegenerate symmetric bilinear form of type~$\epsilon$ (Appendix~\ref{app:type_epsilon}).
\end{lemma}
\begin{proof}
The standard inner product restricted to $x^\perp$ has radical $\langle x \rangle$ (since $(x,x)=0$), inducing a nondegenerate form on the quotient $V := x^\perp / \langle x \rangle$ of dimension $n - 2$. By Lemma~\ref{lem:hyperbolic_completion}, we may choose $v$ such that $\{x, v\}$ forms a hyperbolic plane: $(x, x) = (v, v) = 0$ and $(x, v) = 1$. Since $(x, v) = 1$ implies $x \notin \langle x, v\rangle^\perp$, the subspace $\langle x, v\rangle^\perp$ is a section of the projection $x^\perp \twoheadrightarrow V$, on which the inherited form coincides with the form on~$V$. The orthogonal decomposition $\FF_q^n = \langle x, v\rangle \perp \langle x, v\rangle^\perp$ together with $\disc(\langle x, v\rangle) = -1$ then yields, by multiplicativity of discriminants,
\[
\disc(V) = \disc(\FF_q^n)/\disc(\langle x, v \rangle) = -1 \quad \text{in } \FF_q^*/(\FF_q^*)^2.
\]

A code $C$ with $\Hull(C) = \langle x \rangle$ satisfies $\langle x \rangle \subseteq C \subseteq x^\perp$ (the first inclusion since $x \in C$; the second since $x \in \Hull(C) \subseteq C^\perp$ implies $C \subseteq x^\perp$). Mapping $C \mapsto \bar C := C/\langle x \rangle \subseteq V$ gives a bijection between codes $C$ with $\langle x \rangle \subseteq C \subseteq x^\perp$ and subspaces $\bar C \subseteq V$, with $\dim \bar C = k - 1$.

Under this bijection, $\Hull(\bar C)$ in $V$ equals $\Hull(C)/\langle x \rangle$: for $u \in C$, the class $\bar u \in \bar C$ lies in $\Hull(\bar C)$ if and only if $(u, c) = 0$ for all $c \in C$ (the induced form is well-defined modulo the radical $\langle x \rangle$), which is equivalent to $u \in C^\perp$, and hence to $u \in \Hull(C)$. Therefore, $\Hull(C) = \langle x \rangle$ if and only if $\bar C$ is LCD on $V$. Since $V$ has discriminant $-1$ in $\FF_q^*$, its isometry type is $\epsilon = \chi(-1)$, and the count of such $\bar C$ is $L^{\chi(-1)}(n-2, k-1, q)$.
\end{proof}

\begin{theorem}[Complete Enumeration of GI-Reducible Codes]\label{thm:enumeration_main}
Let $q$ be odd, and let $n, k$ be positive integers with $1 \leq k < n$. The number of GI-reducible $[n,k]_q$ codes is
\[
L(n,k,q) + \frac{K(n,q)}{q-1} \cdot L^{\chi(-1)}(n-2,\, k-1,\, q),
\]
where $L(n,k,q)$ counts standard LCD codes (Appendix~\ref{app:lcd_enumeration}), $K(n,q)$ counts isotropic vectors with nonzero coordinate sum (Lemma~\ref{lem:quadratic_count}), and $L^\epsilon(n,k,q)$ counts LCD subspaces with respect to a type-$\epsilon$ bilinear form (Appendix~\ref{app:type_epsilon}). The exponent $\chi(-1)$ equals $+1$ when $q \equiv 1 \pmod 4$ and $-1$ when $q \equiv 3 \pmod 4$.
\end{theorem}

\begin{proof}
By Theorem~\ref{thm:gi_reducible_count}, GI-reducible codes split into LCD codes (count $L(n,k,q)$) and codes with $\dim(\Hull(C)) = 1$ whose hull vector $x$ satisfies $(x,\1) \neq 0$. The valid hull lines $\langle x \rangle$ correspond to projective points: the set $\{x \in \FF_q^n : (x,x)=0,\, (x,\1) \neq 0\}$ has $K(n,q)$ elements, partitioned into $K(n,q)/(q-1)$ lines (each line contains $q-1$ generators in this set, since the conditions are scalar-invariant). For each line, Lemma~\ref{lem:hull_line_count} gives $L^{\chi(-1)}(n-2, k-1, q)$ codes with that hull.
\end{proof}

\begin{remark}
When $q \equiv 1 \pmod 4$, $\chi(-1) = +1$ and so $L^{\chi(-1)} = L^{+1} = L$, recovering Carlet's standard formula. When $q \equiv 3 \pmod 4$, $\chi(-1) = -1$ and $L^{-1}$ differs from $L$ only in the case $(k\text{ odd}, n\text{ even})$.
\end{remark}

Since codes with $\dim(\Hull(C)) \geq 2$ are never GI-reducible (Theorem~\ref{thm:gi_reducible_count}), this enumeration also provides a criterion for selecting parameters that avoid the orthogonal projector reduction in code-based cryptographic schemes.

\section{Reduction Algorithm}
\label{sec:algorithm}

Having characterized and enumerated GI-reducible codes, we now present the explicit reduction algorithm. The following lemma justifies that, when searching for a structure matrix~$M$ that makes both input codes $M$-LCD, we may restrict attention to a single~$M$ shared by the two codes.

\begin{lemma}[Permutation closure of $M$-LCD codes]
\label{lem:permutation_closure}
Let $M = aI + bJ$ be nondegenerate. If $C$ is $M$-LCD, then $CP$ is $M$-LCD for every permutation matrix~$P$.
\end{lemma}
\begin{proof}
Since $PMP^T = M$, the generator matrix $GP$ of $CP$ satisfies
$$
(GP)M(GP)^T = G(PMP^T)G^T = GMG^T,
$$
so $\rank((GP)M(GP)^T) = \rank(GMG^T) = k$.
\end{proof}

Given two $[n,k]_q$ codes $C_1, C_2$ with generator matrices $G_1, G_2$, the algorithm proceeds as follows:
\begin{enumerate}
\item Check hull dimensions. Compute $\dim(\Hull(C_i))$ for $i=1,2$ (with respect to the standard inner product). If $\dim(\Hull(C_1)) \neq \dim(\Hull(C_2))$, then $C_1 \not\cong C_2$ (since $\Hull(CP) = \Hull(C)P$ for any permutation~$P$, the hull dimension is a permutation invariant). If $\dim(\Hull(C_i)) \geq 2$ for either $i$, then the codes are not GI-reducible by Theorem~\ref{thm:gi_reducible_count}, so return ``not GI-reducible''. If both codes are LCD, set $M = I$ and proceed to Step 3 (Bardet et al.'s construction). Otherwise, proceed to Step 2.

\item Enumerate candidate structure matrices. Since the orthogonal projector $\Pi_{C,M}$ is invariant under scaling $M \mapsto \lambda M$ ($\lambda \in \FF_q^*$), we may fix $a = 1$ and enumerate over $b \in \FF_q^*$ with $1 + nb \neq 0$. For each such $M = I + bJ$, test the LCD condition: compute $\rank(G_i M G_i^T)$ for $i=1,2$. If both equal $k$, then both codes are $M$-LCD; proceed to Step 3. Otherwise, try the next $b$. (In fact, by the proof of Theorem~\ref{thm:gi_reducible_count}, a single valid choice of $b$ suffices whenever both codes have hull dimension $1$ with $(x,\1) \neq 0$.)

\item Construct adjacency matrices. Compute the orthogonal projectors
\[
\Pi_i = M G_i^T (G_i M G_i^T)^{-1} G_i \quad (i=1,2).
\]
Note that $\Pi_i$ is generally not symmetric (unless $M = I$), so $\Pi_i$ represents a directed graph. As noted in Section~\ref{sec:preliminaries}, directed graph isomorphism reduces to GI in polynomial time.

\item Invoke GI solver. Run a graph isomorphism algorithm on $\Pi_1$ and $\Pi_2$. If isomorphic, return ``$C_1 \cong C_2$'' (since $\Pi_{C_1,M} \cong \Pi_{C_2,M}$ implies $C_1 \cong C_2$ by Theorem~\ref{thm:main_characterization}). Otherwise, return ``$C_1 \not\cong C_2$'' (since both codes are $M$-LCD for the same $M$, Theorem~\ref{thm:main_characterization} gives the contrapositive).

\item Exhaustive search. If no $M$ makes both codes $M$-LCD, then $C_1 \not\cong C_2$ (by Lemma~\ref{lem:permutation_closure}).
\end{enumerate}

The procedure is summarized as Algorithm~\ref{alg:gi_reduction}.

\begin{algorithm}[h]
\caption{Reduction from PEP to GI via permutation-invariant bilinear forms}
\label{alg:gi_reduction}
\begin{algorithmic}[1]
\REQUIRE Generator matrices $G_1, G_2 \in \FF_q^{k \times n}$ of $[n,k]_q$ codes $C_1, C_2$.
\ENSURE ``$C_1 \cong C_2$'', ``$C_1 \not\cong C_2$'', or ``not GI-reducible''.
\STATE $h_i \gets k - \rank(G_i G_i^T)$ for $i = 1, 2$ \hfill\COMMENT{hull dimensions}
\IF{$h_1 \neq h_2$}
  \STATE \textbf{return} ``$C_1 \not\cong C_2$''
\ENDIF
\IF{$h_1 \geq 2$}
  \STATE \textbf{return} ``not GI-reducible''
\ENDIF
\IF{$h_1 = 0$}
  \STATE $M \gets I$ \hfill\COMMENT{LCD case: Bardet et al.}
\ELSE
  \STATE find $b \in \FF_q^*$ with $1 + nb \neq 0$ and $\rank\bigl(G_i (I + bJ) G_i^T\bigr) = k$ for $i = 1, 2$
  \IF{no such $b$ exists}
    \STATE \textbf{return} ``$C_1 \not\cong C_2$'' \hfill\COMMENT{by Lemma~\ref{lem:permutation_closure}}
  \ENDIF
  \STATE $M \gets I + bJ$
\ENDIF
\STATE $\Pi_i \gets M G_i^T (G_i M G_i^T)^{-1} G_i$ for $i = 1, 2$
\STATE decide $\Pi_1 \cong \Pi_2$ as edge-weighted directed graphs via a GI solver
\IF{$\Pi_1 \cong \Pi_2$}
  \STATE \textbf{return} ``$C_1 \cong C_2$''
\ELSE
  \STATE \textbf{return} ``$C_1 \not\cong C_2$''
\ENDIF
\end{algorithmic}
\end{algorithm}

Complexity. Step 1 costs $O(k^2 n)$ for computing $GG^T$ and $O(k^3)$ for rank computation. Steps 2--4 enumerate at most $q$ values of $b$. For each value, computing $G_i M G_i^T$ costs $O(k^2 n)$, rank computation and matrix inversion cost $O(k^3)$, and constructing $\Pi_i$ costs $O(kn^2)$. Since $k \leq n$, each iteration costs $O(n^3)$, and invoking the GI solver adds $T_{\mathrm{GI}}(n)$. Summing over at most $q$ structure matrices, the total runtime is $O(q \cdot (n^3 + T_{\mathrm{GI}}(n)))$. For Babai's algorithm~\cite{babai2016}, $T_{\mathrm{GI}}(n) = n^{O(\log^c n)}$ for some constant $c$.

\section{Conclusion}
\label{sec:conclusion}

We established a complete characterization of codes reducible to Graph Isomorphism via the orthogonal projector method of Bardet et al.\ by generalizing the inner product. While the original reduction of Bardet et al.\ applies only to LCD codes using the standard inner product, we proved that replacing the standard inner product with a permutation-invariant bilinear form $M = aI + bJ$ extends the reduction to codes with hull dimension at most~$1$ when $q$ is odd, while in characteristic~$2$ only LCD codes are reducible. This characterization is complete within the orthogonal projector framework: the orthogonal projector gives a valid reduction if and only if the bilinear form has this structure, and codes with hull dimension at least~$2$ cannot be reduced by this method. We derived exact enumeration formulas using Weil's theorem and provided an explicit polynomial-time reduction algorithm.

An open question is whether reduction strategies other than orthogonal projectors can map PEP for codes with $\dim(\Hull(C)) \geq 2$ to GI, or whether such reductions are provably impossible.

\section*{Acknowledgements}
K.~Ishizuka was supported by JSPS KAKENHI Grant Number JP25K17290.

\appendix

\section{LCD Code Enumeration}
\label{app:lcd_enumeration}

This appendix presents the mass formulas for LCD codes, which are central to counting GI-reducible codes. We first review the historical development of LCD code enumeration, then provide explicit formulas separated by field characteristic.

\subsection{Historical Development}

The study of LCD codes was initiated by Massey~\cite{massey1992linear}, who introduced the concept and proved that an $[n,k]$ code $C$ is LCD if and only if its generator matrix $G$ satisfies $\det(GG^T) \neq 0$. Massey's work established the foundational characterization but did not address enumeration.

The first mass formulas for LCD codes were obtained by Carlet et al.~\cite{carlet2019new}, who used the action of the orthogonal group to derive closed-form counting formulas for LCD codes over $\mathbb{F}_2$ and over finite fields of odd characteristic. Their approach exploited the structure of the orthogonal group and its orbits on the Grassmannian of linear codes.

For fields of even characteristic beyond $\mathbb{F}_2$, Li et al.~\cite{lishiling2025} derived mass formulas for linear codes with arbitrary prescribed hull dimension $l$ over finite fields of characteristic 2. In particular, for LCD codes (the case $l=0$), their result extends Carlet's formula from $q=2$ to arbitrary $q=2^m$ by observing that the key structural properties of Carlet's construction generalize naturally to all fields of characteristic 2.

\subsection{Mass Formulas for Odd Characteristic}
\label{app:mass_odd}

Let $q = p^m$ where $p$ is an odd prime. Let $\left[\begin{smallmatrix} n \\ k \end{smallmatrix}\right]_q$ denote the Gaussian binomial coefficient, and let $\chi$ denote the quadratic character of $\mathbb{F}_q$. The number $L(n,k,q)$ of LCD codes of type $[n,k]_q$ is given by Carlet et al.~\cite[Corollary~32]{carlet2019new}:

$$
L(n,k,q) = \begin{cases}
q^{\frac{k(n-k)-1}{2}} \left( q^{\frac{n}{2}} - \chi((-1)^{\frac{n}{2}}) \right) \left[\begin{smallmatrix} \frac{n}{2}-1 \\ \frac{k-1}{2} \end{smallmatrix}\right]_{q^2}, & k \text{ odd, } n \text{ even,} \\[6pt]
q^{\frac{(k+1)(n-k)}{2}} \left[\begin{smallmatrix} \frac{n-1}{2} \\ \frac{k-1}{2} \end{smallmatrix}\right]_{q^2}, & k \text{ odd, } n \text{ odd,} \\[6pt]
q^{\frac{k(n-k+1)}{2}} \left[\begin{smallmatrix} \frac{n-1}{2} \\ \frac{k}{2} \end{smallmatrix}\right]_{q^2}, & k \text{ even, } n \text{ odd,} \\[6pt]
q^{\frac{k(n-k)}{2}} \left[\begin{smallmatrix} \frac{n}{2} \\ \frac{k}{2} \end{smallmatrix}\right]_{q^2}, & k \text{ even, } n \text{ even.}
\end{cases}
$$

\subsection{Mass Formulas for Even Characteristic}

Let $q = 2^m$. Carlet et al.~\cite[Corollary~17]{carlet2019new} derived mass formulas for LCD codes over $\mathbb{F}_2$. Li et al.~\cite[Theorem~2 and subsequent remark]{lishiling2025} showed that Carlet's result for LCD codes over $\mathbb{F}_2$ can be naturally extended to LCD codes over $\mathbb{F}_{2^m}$. The number $L(n,k,q)$ of LCD codes of type $[n,k]_q$ depends on the parity of $n$ and $k$:

$$
L(n,k,q) = \begin{cases}
q^{\frac{nk-k^2+n-1}{2}} \left[\begin{smallmatrix} \frac{n}{2}-1 \\ \frac{k-1}{2} \end{smallmatrix}\right]_{q^2}, & k \text{ odd, } n \text{ even,} \\[6pt]
q^{\frac{(n-k)(k+1)}{2}} \left[\begin{smallmatrix} \frac{n-1}{2} \\ \frac{k-1}{2} \end{smallmatrix}\right]_{q^2}, & k \text{ odd, } n \text{ odd,} \\[6pt]
q^{\frac{k(n-k+1)}{2}} \left[\begin{smallmatrix} \frac{n-1}{2} \\ \frac{k}{2} \end{smallmatrix}\right]_{q^2}, & k \text{ even, } n \text{ odd,} \\[6pt]
q^{\frac{k(n-k)}{2}} \left( q^{n-k} \left[\begin{smallmatrix} \frac{n}{2}-1 \\ \frac{k}{2}-1 \end{smallmatrix}\right]_{q^2} + \left[\begin{smallmatrix} \frac{n}{2}-1 \\ \frac{k}{2} \end{smallmatrix}\right]_{q^2} \right), & k \text{ even, } n \text{ even.}
\end{cases}
$$
\section{Mass Formulas for Type-$\epsilon$ Forms}
\label{app:type_epsilon}

\subsection{Motivation}

By the classification of nondegenerate symmetric bilinear forms over $\FF_q$ ($q$ odd) recalled in Section~\ref{sec:preliminaries}, every such form on $\FF_q^n$ is isometric to exactly one of
\[
M^{+1} := \mathrm{diag}(1, 1, \ldots, 1, 1) \quad \text{or} \quad M^{-1} := \mathrm{diag}(1, 1, \ldots, 1, \gamma),
\]
where $\gamma \in \FF_q^*$ is a fixed non-square. We refer to $\epsilon \in \{+1, -1\}$ as the \emph{type} of the form, and adopt $M^\epsilon$ as the structure matrix throughout this appendix, in line with the $M$-LCD terminology of the main text (Section~\ref{sec:extension}).

Carlet et al.~\cite{carlet2019new} derive their LCD code mass formulas under the implicit assumption that the ambient bilinear form is the standard inner product, i.e., type~$\epsilon = +1$. However, in the Hull-Line Counting Lemma (Lemma~\ref{lem:hull_line_count}), the form induced on the quotient $V = x^\perp / \langle x \rangle$ has discriminant $-1$, which corresponds to type $\epsilon = \chi(-1)$: namely $\epsilon = +1$ when $q \equiv 1 \pmod 4$ and $\epsilon = -1$ when $q \equiv 3 \pmod 4$. To handle the latter case, we generalize Carlet's mass formulas from $\epsilon = +1$ to $\epsilon \in \{+1, -1\}$.

\subsection{Setup and Diagonalisation}

For $\epsilon \in \{+1, -1\}$, with $M^\epsilon$ as above, a code $C$ of length $n$ is \emph{$M^\epsilon$-LCD} (in the sense of Section~\ref{sec:extension}) if $\det(G M^\epsilon G^T) \neq 0$ for some (any) generator matrix~$G$. Define
\[
L^\epsilon(n, k, q) := \bigl|\{C \subseteq \FF_q^n : \dim C = k,\ C \text{ is } M^\epsilon\text{-LCD}\}\bigr|.
\]
The case $\epsilon = +1$ recovers the standard LCD count $L(n, k, q)$ of Appendix~\ref{app:lcd_enumeration}.

A two-fold classification is at play: while $M^\epsilon$ itself has type~$\epsilon$, each $M^\epsilon$-LCD code~$C$ has its own \emph{inner discriminant class} $\delta_C \in \{1, \gamma\}$, defined by the property that some generator matrix~$G$ satisfies $G M^\epsilon G^T = \mathrm{diag}(1, \ldots, 1, \delta_C)$ (existence by~\cite[Theorem~25]{carlet2019new}). Correspondingly, we partition
\[
L^\epsilon(n, k, q) = |\mathrm{LCD}^\epsilon_+[n, k]_q| + |\mathrm{LCD}^\epsilon_-[n, k]_q|,
\]
where the subscript indicates $\delta_C = 1$ or $\delta_C = \gamma$, respectively.

\subsection{Reference Codes and Orbit Decomposition}

Let $1 \leq k < n$ and assume $M^\epsilon = \mathrm{diag}(1, \ldots, 1, \tau_\epsilon)$ with $\tau_{+1} := 1$, $\tau_{-1} := \gamma$. Denote the standard basis of $\FF_q^n$ by $e_1, \ldots, e_n$. Recall (\cite[Lemma~26]{carlet2019new}; also a direct counting argument since the squares form a subgroup of index $2$ in $\FF_q^*$ for $q$ odd) that every element of $\FF_q$ is a sum of two squares; fix $a, b \in \FF_q$ with $a^2 + b^2 = \gamma$.

\begin{lemma}[Reference codes]\label{lem:reference_codes}
Let $1 \leq k < n$ and $\epsilon \in \{+1, -1\}$, with $M^\epsilon$ as above. Then:
\begin{enumerate}
\item[(i)] The code $\mathcal{C}_+^\epsilon := \langle e_1, \ldots, e_k\rangle$, generated by $G_+^\epsilon := (I_k \mid 0_{k \times (n-k)})$, is $M^\epsilon$-LCD with inner discriminant $\delta_+ = 1$.
\item[(ii)] Define $v \in \FF_q^n$ by
\begin{itemize}
\item $v := a\, e_k + b\, e_{k+1}$ if $k \leq n-2$;
\item $v := a\, e_{n-1} + b\, e_n$ if $k = n-1$ and $\epsilon = +1$;
\item $v := e_n$ if $k = n-1$ and $\epsilon = -1$.
\end{itemize}
The code $\mathcal{C}_-^\epsilon := \langle e_1, \ldots, e_{k-1}, v\rangle$, generated by the matrix $G_-^\epsilon$ with rows $e_1, \ldots, e_{k-1}, v$, is $M^\epsilon$-LCD with inner discriminant $\delta_- = \gamma$.
\end{enumerate}
\end{lemma}
\begin{proof}
For (i), since $k \leq n-1$, none of the rows of $G_+^\epsilon$ touches the last column of $M^\epsilon$, so $G_+^\epsilon M^\epsilon (G_+^\epsilon)^T = I_k$, which has determinant $1 \neq 0$. Hence $\mathcal{C}_+^\epsilon$ is $M^\epsilon$-LCD with $\delta_+ = 1$.

For (ii), $v$ is orthogonal (under $M^\epsilon$) to $e_1, \ldots, e_{k-1}$ by construction. The value $v M^\epsilon v^T$ equals $a^2 + b^2 = \gamma$ in the first two cases, and $\tau_{-1} = \gamma$ in the third. Hence $G_-^\epsilon M^\epsilon (G_-^\epsilon)^T = \mathrm{diag}(1, \ldots, 1, \gamma)$, which has determinant $\gamma \neq 0$, so $\mathcal{C}_-^\epsilon$ is $M^\epsilon$-LCD with $\delta_- = \gamma$.
\end{proof}

For $\delta \in \FF_q^*$, define the orthogonal group
\[
\mathbb{O}_m^\delta(q) := \{Q \in \mathrm{GL}_m(\FF_q) : Q \mathrm{diag}(1, \ldots, 1, \delta) Q^T = \mathrm{diag}(1, \ldots, 1, \delta)\}.
\]
The parameter $\delta$ serves as a generic discriminant placeholder; below it will be instantiated as $1$, $\gamma$, $\epsilon$, or $\epsilon\gamma$ depending on context. In particular, $\mathbb{O}_n^\epsilon(q)$ is the orthogonal group preserving $M^\epsilon$, and its transitivity on each $\mathrm{LCD}^\epsilon_\pm[n, k]_q$ follows directly from Witt's extension theorem.

\begin{lemma}[Transitivity via Witt extension]\label{lem:witt_transitivity}
For $\epsilon \in \{+1, -1\}$ and a fixed inner discriminant class $\delta \in \{1, \gamma\}$, the group $\mathbb{O}_n^\epsilon(q)$ acts transitively on the set of $M^\epsilon$-LCD codes of dimension $k$ with inner discriminant $\delta$.
\end{lemma}
\begin{proof}
Let $C_1, C_2$ be two such codes. Each $C_i$ is equipped with the nondegenerate symmetric bilinear form $M^\epsilon|_{C_i}$ on a $k$-dimensional space, of common discriminant class $\delta$. By the classification of nondegenerate symmetric bilinear forms over $\FF_q$ ($q$ odd, Section~\ref{sec:preliminaries}), two such forms on a fixed-dimensional space are isometric if and only if they have equal discriminant class; hence there exists a linear isomorphism $\phi: C_1 \to C_2$ with $\phi(x) M^\epsilon \phi(y)^T = x M^\epsilon y^T$ for all $x, y \in C_1$. By Witt's extension theorem (Theorem~\ref{thm:witt_extension}) applied to $(\FF_q^n, M^\epsilon)$, $\phi$ extends to an element $\Phi \in \mathbb{O}_n^\epsilon(q)$, and $\Phi(C_1) = C_2$.
\end{proof}

By Lemma~\ref{lem:witt_transitivity} applied to the reference codes $\mathcal{C}_\pm^\epsilon$ of Lemma~\ref{lem:reference_codes}, $\mathrm{LCD}^\epsilon_\pm[n, k]_q = \mathcal{C}_\pm^\epsilon \cdot \mathbb{O}_n^\epsilon(q)$, so by orbit--stabilizer,
\begin{equation}\label{eq:orbit_stab}
|\mathrm{LCD}^\epsilon_\pm[n, k]_q| = \frac{|\mathbb{O}_n^\epsilon(q)|}{|\mathrm{St}(\mathcal{C}_\pm^\epsilon)|}.
\end{equation}

\subsection{Stabilisers and Orthogonal Group Orders}

By~\cite[Corollary~30]{carlet2019new}, the stabiliser of $\mathcal{C}_\pm^\epsilon$ is block-diagonal in the basis adapted to $\mathcal{C}_\pm^\epsilon \perp_{M^\epsilon} (\mathcal{C}_\pm^\epsilon)^{\perp_{M^\epsilon}}$:
\begin{align*}
|\mathrm{St}(\mathcal{C}_+^\epsilon)| &= |\mathbb{O}_k^1(q)| \cdot |\mathbb{O}_{n-k}^{\epsilon}(q)|, \\
|\mathrm{St}(\mathcal{C}_-^\epsilon)| &= |\mathbb{O}_k^\gamma(q)| \cdot |\mathbb{O}_{n-k}^{\epsilon\gamma}(q)|,
\end{align*}
where the bottom-block discriminants come from the multiplicativity $\delta_C \cdot \delta_{C^\perp} \equiv \epsilon \pmod{(\FF_q^*)^2}$. The order of $\mathbb{O}_m^\delta(q)$ is given by~\cite[Eqs.~(11)--(12)]{carlet2019new} (citing Weyl~\cite{weyl1939classical}): for $n$ odd,
\[
|\mathbb{O}_n^\delta(q)| = 2 q^{(n-1)^2/4} \prod_{i=1}^{(n-1)/2}(q^{2i} - 1) \quad (\text{independent of } \delta),
\]
and for $n$ even,
\[
|\mathbb{O}_n^\delta(q)| = 2 q^{n(n-2)/4} \bigl(q^{n/2} - \chi((-1)^{n/2}\delta)\bigr) \prod_{i=1}^{n/2 - 1}(q^{2i} - 1).
\]

\subsection{Mass Formula}

Substituting orders into~\eqref{eq:orbit_stab},
\begin{equation}\label{eq:LB_pm}
\begin{aligned}
|\mathrm{LCD}^\epsilon_+[n,k]_q| &= \frac{|\mathbb{O}_n^\epsilon(q)|}{|\mathbb{O}_k^1(q)| \cdot |\mathbb{O}_{n-k}^{\epsilon}(q)|}, \\
|\mathrm{LCD}^\epsilon_-[n,k]_q| &= \frac{|\mathbb{O}_n^\epsilon(q)|}{|\mathbb{O}_k^\gamma(q)| \cdot |\mathbb{O}_{n-k}^{\epsilon\gamma}(q)|}.
\end{aligned}
\end{equation}

\begin{theorem}[Mass formula for $L^\epsilon$]\label{thm:Lepsilon}
For $\epsilon \in \{+1, -1\}$, $q$ odd, and $1 \leq k < n$,
\[
L^\epsilon(n, k, q) = \begin{cases}
q^{\frac{k(n-k)-1}{2}}\bigl(q^{n/2} - \chi((-1)^{n/2}\epsilon)\bigr)\left[\begin{smallmatrix} \frac{n}{2}-1 \\ \frac{k-1}{2} \end{smallmatrix}\right]_{q^2}, & k \text{ odd, } n \text{ even,} \\[6pt]
L(n, k, q), & \text{otherwise.}
\end{cases}
\]
In particular, $L^{+1}(n, k, q) = L(n, k, q)$ recovers Carlet's formula~\cite[Corollary~32]{carlet2019new}.
\end{theorem}

\begin{proof}
We compute $L^\epsilon = |\mathrm{LCD}^\epsilon_+| + |\mathrm{LCD}^\epsilon_-|$ from~\eqref{eq:LB_pm} and compare with $L = L^{+1}$, the case treated by Carlet et al.~\cite[Corollary~32]{carlet2019new}.

When $m$ is even, write $|\mathbb{O}_m^\delta| = P_m \cdot A_m^\delta$, where $P_m := 2 q^{m(m-2)/4} \prod_{i=1}^{m/2 - 1}(q^{2i} - 1)$ is $\delta$-independent and $A_m^\delta := q^{m/2} - \chi((-1)^{m/2}\delta)$. Using $\chi(\gamma) = -1$ and $\chi((-1)^{m/2}\epsilon)^2 = 1$, the pair $\{A_m^\epsilon, A_m^{\epsilon\gamma}\}$ satisfies
\[
A_m^\epsilon + A_m^{\epsilon\gamma} = 2 q^{m/2},
\qquad
A_m^\epsilon \cdot A_m^{\epsilon\gamma} = q^m - 1,
\]
both $\epsilon$-independent. When $m$ is odd, $|\mathbb{O}_m^\delta|$ does not depend on $\delta$. We treat the four parity cases.

\emph{Case $(k\text{ odd}, n\text{ odd})$.} Both $|\mathbb{O}_n^\delta|$ and $|\mathbb{O}_k^\delta|$ are $\delta$-independent, so
\[
L^\epsilon
= \frac{|\mathbb{O}_n|}{|\mathbb{O}_k|}\Bigl[\frac{1}{|\mathbb{O}_{n-k}^\epsilon|} + \frac{1}{|\mathbb{O}_{n-k}^{\epsilon\gamma}|}\Bigr]
= \frac{|\mathbb{O}_n|}{|\mathbb{O}_k|} \cdot \frac{A_{n-k}^\epsilon + A_{n-k}^{\epsilon\gamma}}{P_{n-k}\, A_{n-k}^\epsilon A_{n-k}^{\epsilon\gamma}},
\]
which is $\epsilon$-independent by the relations above.

\emph{Case $(k\text{ even}, n\text{ odd})$.} $|\mathbb{O}_n^\delta|$ and $|\mathbb{O}_{n-k}^\delta|$ are $\delta$-independent; the only $\delta$-dependence is in $|\mathbb{O}_k^\delta|$, which appears as $|\mathbb{O}_k^1|$ in $|\mathrm{LCD}^\epsilon_+|$ and $|\mathbb{O}_k^\gamma|$ in $|\mathrm{LCD}^\epsilon_-|$ irrespective of $\epsilon$. Hence $L^\epsilon$ has no $\epsilon$-dependence.

\emph{Case $(k\text{ even}, n\text{ even})$.} All three orders depend on $\delta$. Set $X = q^{k/2}$, $Y = q^{(n-k)/2}$, $s = \chi((-1)^{k/2})$, $t = \chi((-1)^{(n-k)/2})$, so that $A_n^\epsilon = XY - \chi(\epsilon)\, st$, $A_k^\delta = X - \chi(\delta) s$, $A_{n-k}^\delta = Y - \chi(\delta) t$. Direct computation gives
\[
\begin{aligned}
L^\epsilon
&= \frac{P_n A_n^\epsilon}{P_k P_{n-k}}\Bigl[\frac{1}{A_k^1 A_{n-k}^\epsilon} + \frac{1}{A_k^\gamma A_{n-k}^{\epsilon\gamma}}\Bigr]
= \frac{P_n A_n^\epsilon}{P_k P_{n-k}} \cdot \frac{2(XY + \chi(\epsilon)\, st)}{(X^2 - 1)(Y^2 - 1)} \\
&= \frac{P_n}{P_k P_{n-k}} \cdot \frac{2((XY)^2 - 1)}{(X^2 - 1)(Y^2 - 1)},
\end{aligned}
\]
using $\chi(\epsilon)^2 = 1$ and $s^2 = t^2 = 1$. Hence $L^\epsilon$ is $\epsilon$-independent.

\emph{Case $(k\text{ odd}, n\text{ even})$.} $|\mathbb{O}_k^\delta|$ and $|\mathbb{O}_{n-k}^\delta|$ are both $\delta$-independent, so $|\mathrm{LCD}^\epsilon_+| = |\mathrm{LCD}^\epsilon_-|$ and
\[
L^\epsilon(n, k, q) = \frac{2|\mathbb{O}_n^\epsilon|}{|\mathbb{O}_k| \cdot |\mathbb{O}_{n-k}|}.
\]
The ratio $L^\epsilon/L = A_n^\epsilon/A_n^{+1} = (q^{n/2} - \chi((-1)^{n/2}\epsilon))/(q^{n/2} - \chi((-1)^{n/2}))$ substituted into Carlet's formula (Appendix~\ref{app:mass_odd}) gives the stated expression.
\end{proof}

In particular, $L^{-1} = L$ when $\chi(-1) = +1$ (i.e., $q \equiv 1 \pmod 4$, where $\chi((-1)^{n/2}\cdot(-1)) = \chi((-1)^{n/2})$); when $\chi(-1) = -1$, the $\chi$-correction flips sign relative to $L$ in the $(k\text{ odd}, n\text{ even})$ case.

\bibliographystyle{plain}
\bibliography{main}

\end{document}
